\section{Measurable functions and the integral}

\emph{Lecture 2, 14 September.}

\begin{definition}
A function $f \colon X \to \R$ is \emph{measurable} if $f^{-1}(B) \in \cF$ for every Borel set $B \subseteq \R$. It is enough to check the sets $\{ f > a \}$ for $a \in \R$.
\end{definition}

Sums, products, pointwise limits, suprema, and infima of measurable functions are measurable.

\begin{definition}
A \emph{simple function} is a finite sum $s = \sum_{i=1}^k a_i \mathbf{1}_{A_i}$ with $a_i \ge 0$ and $A_i \in \cF$. Its integral is $\int s \, d\mu = \sum_{i=1}^k a_i \mu(A_i)$. For measurable $f \ge 0$ we define
\[
  \int f \, d\mu = \sup \Bigl\{ \int s \, d\mu : 0 \le s \le f, \ s \text{ simple} \Bigr\}.
\]
\end{definition}

\begin{theorem}[Monotone convergence]
\label{thm:mct}
If $0 \le f_1 \le f_2 \le \cdots$ are measurable and $f_n \to f$ pointwise, then
\[
  \lim_{n \to \infty} \int f_n \, d\mu = \int f \, d\mu .
\]
\end{theorem}

\begin{remark}
Monotonicity cannot be dropped. On $[0, 1]$ with Lebesgue measure, $f_n = n \mathbf{1}_{(0, 1/n)}$ converges to $0$ pointwise, but $\int f_n = 1$ for every $n$.
\end{remark}
